\src{190}
variant dans la catégorie des faisceaux algébriques cohérents sur $\symup{X}$, et dont
les valeurs sont dans la catégorie des faisceaux algébriques cohérents sur $\symup{X}$,
resp. dans la catégorie des modules sur $\symup{H^{0}(X,\Osh_X)}$, par les formules :
\begin{equation}\tag{8.1}
\resizebox{\textwidth}{!}{$
\left\{
\begin{aligned}
\symup{\underline{T}_{r}^{s*}(A_1,\ldots,A_r;B_1,\ldots,B_s)}
&
\symup{={}H^{*}\!\left(
    \Homsh_{\Osh_X}\!\left(
      \underline{L}(A_1)\otimes\cdots\otimes \underline{L}(A_r),
      \underline{L}(B_1)\otimes\cdots\otimes \underline{L}(B_s)
    \right)\right)}
\\[3pt]
\symup{T_{r}^{s*}(A_1,\ldots,A_r;B_1,\ldots,B_s)}
&
\symup{={}\underline{\underline{R}}^{*}\Gamma\!\left(
    \Homsh_{\Osh_X}\!\left(
      \underline{L}(A_1)\otimes\cdots\otimes \underline{L}(A_r),
      \underline{L}(B_1)\otimes\cdots\otimes \underline{L}(B_s)
    \right)\right)}
\end{aligned}
\right.
$}
\end{equation}
Dans cette formule, $\symup{\underline{L}(F)}$ désigne une résolution finie localement libre du
faisceau algébrique cohérent $\symup{F}$, et $\symup{\underline{\underline{R}}^{*}\Gamma(K)}$ désigne
l'hypercohomologie de l'espace $\symup{X}$ par rapport au complexe de faisceaux $\symup{K}$. Si
$\symup{r}$ ou $\symup{s}$ est nul, on remplace le produit tensoriel des $\symup{\underline{L}(A_i)}$, resp. des
$\symup{\underline{L}(B_j)}$, par $\symup{\Osh_X}$. En particulier $\symup{\underline{T}_{0}^{0}}$ et $\symup{T_{0}^{0}}$
sont des foncteurs gradués à 0 arguments, $\symup{\underline{T}_{0}^{0}}$ se réduit au
degré 0, où il est le faisceau $\symup{\Osh_X}$, tandis que $\symup{T_{0}^{0}}$ est identique à
$\symup{H^{*}(X,\Osh_X)}$. Le fait que les deuxièmes membres de (8.1) ne dépendent pas
des résolutions choisies est de toutes façons évident pour la première ligne
(la question étant alors locale), et pour la deuxième résulte des remarques
générales qui précédaient, compte tenu de la suite spectrale
d'hypercohomologie du complexe de faisceaux
$\symup{K=\Homsh_{\Osh_X}\bigl(\underline{L}(A_1)\otimes\cdots,
      \underline{L}(B_1)\otimes\cdots\bigr)}$, aboutissant à l'hypercohomologie de $\symup{X}$ par rapport à $\symup{K}$, et dont le terme
initial est $\symup{H^{p}(X,H^{q}(K))}$, i.e.
\begin{equation}\tag{8.2}
  \symup{E_{2}^{p,q}
  =H^{p}\!\left(X,
    (\underline{T}_{r}^{s})^{(q)}
      (A_1,\ldots,A_r;B_1,\ldots,B_s)\right)}
\end{equation}
On voit alors que cette suite spectrale elle-même ne dépend pas des résolutions
choisies, son aboutissement est le premier membre de (8.1). On définit
facilement les opérateurs cobords relatifs aux divers arguments $\symup{A_i}$, $\symup{B_j}$,
en notant que toute suite exacte
$\symup{0\longrightarrow F'\longrightarrow F\longrightarrow F''\longrightarrow0}$
peut se résoudre par une suite exacte de complexes finis localement libres.

On définit dans le système des foncteurs $\symup{\underline{T}_{r}^{s*}}$, resp.
$\symup{T_{r}^{s*}}$, des opérations analogues à celles du calcul tensoriel, et dont la
définition est immédiate à partir des formules de définition (8.1). Ainsi, on a
une \uline{composition} (généralisant celle envisagée dans le n\textsuperscript{o} 2) :
\begin{equation}\tag{8.3}
\resizebox{\textwidth}{!}{$\symup{
T_{r}^{s*}(A_1,\ldots,A_r;B_1,\ldots,B_s)
\times T_{r'}^{s'*}(A'_1,\ldots,A'_{r'};B'_1,\ldots,B'_{s'})
\longrightarrow T_{r+r'}^{s+s'*}
(A_1,\ldots,A'_{r'};B_1,\ldots,B'_{s'})
}$}
\end{equation}
satisfaisant aux propriétés évidentes d'associativité, de compatibilité avec
les homomorphismes fonctoriels et les homomorphismes cobords, les suites
spectrales. De même, on a des opérations de symétrie qu'on laisse au lecteur le
soin d'expliciter. On a de plus une opération de \uline{contraction} chaque fois que
l'un des arguments $\symup{A_i}$

